\begin{table}[!h]
\centering
\caption{\label{tab:raiz}Testes de raiz unitária (níveis em log).}
\centering
\fontsize{10}{12}\selectfont
\begin{tabular}[t]{lrrrr}
\toprule
\textbf{Região} & \textbf{ADF nível} & \textbf{ADF 1\textsuperscript{a} dif.} & \textbf{PP nível} & \textbf{KPSS nível}\\
\midrule
\cellcolor{gray!10}{Média Brasil} & \cellcolor{gray!10}{-3,968} & \cellcolor{gray!10}{-6,378} & \cellcolor{gray!10}{-2,746} & \cellcolor{gray!10}{0,276}\\
Bahia & -3,290 & -5,913 & -2,398 & 0,293\\
\cellcolor{gray!10}{Espírito Santo} & \cellcolor{gray!10}{-3,506} & \cellcolor{gray!10}{-5,994} & \cellcolor{gray!10}{-2,651} & \cellcolor{gray!10}{0,227}\\
Goiás & -3,798 & -6,694 & -2,686 & 0,292\\
\cellcolor{gray!10}{Minas Gerais} & \cellcolor{gray!10}{-4,116} & \cellcolor{gray!10}{-6,385} & \cellcolor{gray!10}{-2,833} & \cellcolor{gray!10}{0,253}\\
Paraná & -4,067 & -6,207 & -2,937 & 0,257\\
\cellcolor{gray!10}{Rio de Janeiro} & \cellcolor{gray!10}{-2,837} & \cellcolor{gray!10}{-5,425} & \cellcolor{gray!10}{-2,494} & \cellcolor{gray!10}{0,331}\\
Santa Catarina & -3,759 & -6,407 & -2,828 & 0,278\\
\cellcolor{gray!10}{São Paulo} & \cellcolor{gray!10}{-4,062} & \cellcolor{gray!10}{-6,223} & \cellcolor{gray!10}{-2,639} & \cellcolor{gray!10}{0,290}\\
\bottomrule
\end{tabular}
\end{table}
